\documentclass[10pt,a4paper]{amsart}
\usepackage[margin=19mm]{geometry}
\usepackage{amsmath,amssymb,mathtools}
\usepackage[scr=rsfs,cal=euler]{mathalfa}
\usepackage{xfrac}
\usepackage[unicode,hidelinks,bookmarks=false]{hyperref}
\newcommand{\ie}{i.e.\ }
\newcommand{\cf}{cf.\ }
\newcommand{\resp}{respectively\ }
\newcommand{\Subsection}{\subsection}
\providecommand{\nobreakdash}{\nobreak\hbox{-}}
\setlength{\emergencystretch}{2em}
\multlinegap0pt
\usepackage{tikz}
\usetikzlibrary{cd}
\usepackage{mathtools}
\newtheorem{lemma}{Lemma}
\renewcommand{\ge}{\geqslant}
\title{Hodge-to-de Rham degeneration and quasihomogeneous singularities of curves}
\author{Yunfan He}
\date{Source: arXiv:2604.05825v1 (CC BY 4.0)}
\begin{document}
\maketitle

\noindent\emph{Source and licence.} Hodge-to-de Rham degeneration and quasihomogeneous singularities of curves. Version arXiv:2604.05825v1 (CC BY 4.0), distributed by arXiv under the Creative Commons Attribution 4.0 International licence, http://creativecommons.org/licenses/by/4.0/, as recorded in the arXiv licence metadata for this exact version and shown on the abstract page https://arxiv.org/abs/2604.05825v1. Full licence notice: \texttt{licenses/arxiv-2604.05825-CC-BY-4.0.txt}.\par\smallskip

\noindent\textit{Editorial scope.} Counted unit: the article's lemma lem:tail-map-euler (source0.tex lines 550--564). The article's complete original proof of it is delivered with it (source0.tex lines 566--742, one \texttt{proof} environment of 177 lines). No mathematics is written, repaired or re-derived: the statement and the proof are the article's own lines, in the article's own notation, and no retained line is substituted or re-flowed.  Retained uncounted as the source-local context the proof needs: lines 89--91; lines 425--448.  Nothing was omitted from any retained span, so the package carries no omission record.  No retained line contains a citation, so the wrapper carries no bibliography and no bibliography entry.  The retained lines contain the article's own diagram code, which the wrapper typesets with the standard tikz packages; no image file is delivered and the package declares no asset dependency.  Editorial formatting only, in addition: an AMS \texttt{amsart} preamble that restates the article's own declarations and macros used by the retained lines, namely the environments \texttt{lemma} on separate counters, together with the standard packages those lines need.  The retained environments print in this document as Lemma 1 in order of appearance, whereas the article prints the same items with the numbering of its own full document.  The complete native source of the article as distributed in the arXiv e-print for this exact version is bundled unchanged as \texttt{sources/197982/source0.tex}.
\begin{equation}\label{eq:weighted-euler}
f=w_1u f_u+w_2v f_v
\end{equation}

\begin{lemma}\label{lem:tail-map-quotient}
Let
\[
  R=\mathbb C[[u,v]],\qquad A=R/(f),
\]
where $f\in R$ defines an isolated plane curve singularity. We use
\[
  (0:_{M_f}f):=\{m\in M_f\mid fm=0\}
\]
to denote the annihilator of \(f\) in \(M_f\). Then for every $p\ge 1$ there are canonical isomorphisms
\[
H^{-p}\!\left(\bigwedge^{p+1}\mathbb{L}^{\bullet}_A\right)\cong (0:_{M_f}f),
\qquad
H^{-p}\!\left(\bigwedge^{p+2}\mathbb{L}^{\bullet}_A\right)\cong T_f,
\]
In particular, their dimensions both equal to the Tjurina number:
\[
  \dim_{\mathbb C}H^{-p}\!\left(\bigwedge^{p+1}\mathbb{L}^{\bullet}_A\right)
  =
  \dim_{\mathbb C}H^{-p}\!\left(\bigwedge^{p+2}\mathbb{L}^{\bullet}_A\right)
  =
  \tau_x.
\]
\end{lemma}

\begin{lemma}\label{lem:tail-map-euler}
Keep the setting of Lemma~\ref{lem:tail-map-quotient}, and assume further that \(f=w_1u f_u+w_2v f_v\) is weighted homogeneous. Then for every $p\ge 1$, the cohomology groups of Lemma~\ref{lem:tail-map-quotient} are both isomorphic to \(M_f\):
\[
H^{-p}\!\left(\bigwedge^{p+1}\mathbb{L}^{\bullet}_A\right)\cong M_f,
\qquad
H^{-p}\!\left(\bigwedge^{p+2}\mathbb{L}^{\bullet}_A\right)\cong M_f.
\]
Moreover, under these identifications, the de Rham differential induces an isomorphism
\[
d_1:
H^{-p}\!\left(\bigwedge^{p+1}\mathbb{L}^{\bullet}_A\right)
\longrightarrow
H^{-p}\!\left(\bigwedge^{p+2}\mathbb{L}^{\bullet}_A\right).
\]
\end{lemma}

\begin{proof}
By Lemma~\ref{lem:tail-map-quotient},
\[
H^{-p}\!\left(\bigwedge^{p+1}\mathbb{L}^{\bullet}_A\right)\cong (0:_{M_f}f),
\qquad
H^{-p}\!\left(\bigwedge^{p+2}\mathbb{L}^{\bullet}_A\right)\cong T_f.
\]
Since Euler relation~\eqref{eq:weighted-euler} gives \(f\in (f_u,f_v)\), multiplication by \(f\) on \(M_f\) is zero. Therefore
\[
(0:_{M_f}f)=M_f,
\qquad
T_f=M_f/fM_f\cong M_f,
\]
which gives the two displayed identifications in the statement.

To compute the de Rham differential, we use the explicit three-term complex from the proof of Lemma~\ref{lem:tail-map-quotient}:
\[
K:
A \xrightarrow{\Delta_2} A^{\oplus 2}\xrightarrow{\Delta_1} A,
\]
where
\[
\Delta_2(c)=(cf_u,cf_v),\qquad \Delta_1(a,b)=-af_v+bf_u.
\]
There we identified
\[
H^{-p}\!\left(\bigwedge^{p+1}\mathbb{L}^{\bullet}_A\right)=H^{-1}(K),
\qquad
H^{-p}\!\left(\bigwedge^{p+2}\mathbb{L}^{\bullet}_A\right)=H^{0}(K).
\]
By Lemma~\ref{lem:tail-map-quotient} and the Euler relation, these two groups are already abstractly isomorphic to \(M_f\). We now construct an explicit identification \(M_f\cong H^{-1}(K)\), which will allow us to compute the de Rham differential on the chain level.

Next, consider the Euler syzygy
\[
E:=(-w_2v,w_1u)\in A^{\oplus 2}.
\]
Because
\[
\Delta_1(E)=w_2v f_v+w_1u f_u=f=0\quad\text{in }A,
\]
the element $E$ defines a class in $H^{-1}(K)$. Multiplication by $E$ yields a map
\[
\theta:M_f\longrightarrow H^{-1}(K),\qquad [h]\longmapsto [hE].
\]
This map is well-defined, since in \(A^{\oplus 2}\) one has
\[
f_uE=\Delta_2(-w_2v),
\qquad
f_vE=\Delta_2(w_1u).
\]
We claim that \(\theta\) is surjective. Let \([(a,b)]\in H^{-1}(K)\), so
\[
-af_v+bf_u=0
\qquad\text{in }A.
\]
Choose lifts \(\widetilde a,\widetilde b\in R\). Then
\[
-\widetilde a f_v+\widetilde b f_u\in (f),
\]
so there exists \(\widetilde c\in R\) such that
\[
-\widetilde a f_v+\widetilde b f_u=\widetilde c\,f.
\]
Using the Euler relation~\eqref{eq:weighted-euler}, we get
\[
(\widetilde b-\widetilde c\,w_1u)f_u+(-\widetilde a-\widetilde c\,w_2v)f_v=0
\qquad\text{in }R.
\]
As noted in the proof of Lemma~\ref{lem:tail-map-quotient}, \((f_u,f_v)\) is a regular sequence in \(R\). Therefore there exists \(\widetilde d\in R\) such that
\[
\widetilde b-\widetilde c\,w_1u=-\widetilde d\,f_v,
\qquad
-\widetilde a-\widetilde c\,w_2v=\widetilde d\,f_u.
\]
After modulo \((f)\), we obtain
\[
b-\overline{\widetilde c}\,w_1u=-\overline{\widetilde d}\,f_v,
\qquad
-a-\overline{\widetilde c}\,w_2v=\overline{\widetilde d}\,f_u
\qquad
\text{in }
A.
\]
Equivalently,
\[
(a,b)=\overline{\widetilde c}\,E+\Delta_2(-\overline{\widetilde d}),
\]
so \([(a,b)]=\theta([\overline{\widetilde c}])\). Hence \(\theta\) is surjective.

Since \(M_f\) is finite-dimensional and \(H^{-1}(K)\cong (0:_{M_f}f)=M_f\), the surjective map \(\theta:M_f\to H^{-1}(K)\) is an isomorphism, and
\[
H^{-1}(K)\cong M_f.
\]

We now compute the de Rham differential on this explicit model. Since
\[
d_{\mathrm{dR}}(d\varepsilon)=0,
\]
one has \(d_{\mathrm{dR}}(\gamma_m)=0\) for every \(m\ge 0\). Therefore the ordinary de Rham differential induces a chain map
\[
\delta_p:K[-p+1]\longrightarrow K[-p]
\]
whose nontrivial components are given by
\[
\begin{tikzcd}[column sep=large, row sep=large]
A
  \arrow[r, "\Delta_2"]
  \arrow[d, "\delta_p^{-p-1}" left]
&
A^{\oplus 2}
  \arrow[r, "\Delta_1"]
  \arrow[d, "\delta_p^{-p}"]
&
A
  \arrow[d, "0"]
\\
A^{\oplus 2}
  \arrow[r, "\Delta_1"]
&
A
  \arrow[r, "0"]
&
0,
\end{tikzcd}
\]
where
\[
\delta_p^{-p-1}(c):=(c_u,c_v),
\]
\[
\delta_p^{-p}(a,b):=b_u-a_v.
\]
Indeed,
\[
d_{\mathrm{dR}}(c)=c_u\,du+c_v\,dv
\]
gives \(\delta_p^{-p-1}(c)=(c_u,c_v)\), while
\[
d_{\mathrm{dR}}(a\,du+b\,dv)=(b_u-a_v)\,du\wedge dv.
\]
Under the identification \(H^{-1}(K)\cong M_f\) via \(\theta\), the induced map on cohomology is computed by
\[
[h]\longmapsto [\delta_p^{-p}(hE)].
\]
Since
\[
hE=(-h w_2v,\; h w_1u),
\]
we obtain
\[
\delta_p^{-p}(hE)=\partial_u(hw_1u)-\partial_v(-hw_2v)
=
w_1(uh_u+h)+w_2(vh_v+h).
\]
Thus
\[
d_1([h])
=
\bigl[w_1(uh_u+h)+w_2(vh_v+h)\bigr]
\in M_f.
\]

If $[h]$ is weighted homogeneous of weighted degree $\lambda$, then Euler's formula gives
\[
w_1uh_u+w_2vh_v=\lambda h
\quad\text{in }M_f,
\]
hence
\begin{equation}\label{eq:weighted-tail-scalar}
d_1([h])=(\lambda+w_1+w_2)[h].
\end{equation}
Since $f_u$ and $f_v$ are weighted homogeneous, the Jacobian ideal $(f_u,f_v)$ is graded for the induced weighted grading on $R$, so
\[
M_f=\bigoplus_{\lambda}(M_f)_\lambda
\]
is a finite direct sum of weighted-homogeneous pieces. By \eqref{eq:weighted-tail-scalar}, $d_1$ acts on $(M_f)_\lambda$ by the scalar \(\lambda+w_1+w_2\), which is nonzero because \(w_1,w_2>0\) and \(\lambda\ge 0\). Therefore $d_1$ is an isomorphism.
\end{proof}

\end{document}
